\begin{textsample}{POS Dim 3 – ce – Score 31.00 – ce/ce\_inf\_13.txt}  \label{ex:f3_pos_052}
4.1.4.
Parceria famílias e \textbf{instituição} de Educação \textbf{Infantil} Educar e \textbf{cuidar} dos \textbf{bebês}, \textbf{crianças} bem \textbf{pequenas} e \textbf{crianças} \textbf{pequenas} requer uma grande parceria e respeito entre escola e família.
Ambas as \textbf{instituições} ( famílias e escola ) têm muitos saberes e informações a compartilhar em relação aos \textbf{cuidados} e aprendizagens das \textbf{crianças}.
Quando uma família matricula sua \textbf{criança} em \textbf{creches} e \textbf{pré-escolas}, precisa compartilhar responsabilidades, deveres e direitos com todas/todos as/os profissionais que lidam com a \textbf{criança} dentro da \textbf{instituição}, assumindo de forma corresponsável os papéis de educar e \textbf{cuidar}.
Professoras/professores e profissionais firmam a parceria e estabelecem vínculos de confiança e de relações afetivas com as \textbf{crianças} e suas famílias.
A parceria entre \textbf{instituição} e família favorece na compreensão do contexto familiar em que a \textbf{criança} está inserida, possibilitando o compartilhamento dos saberes e o respeito a sua cultura.
Assim como as concepções de \textbf{criança}, \textbf{infância} e Educação \textbf{Infantil} foram sendo transformadas ao longo dos tempos, comprovados com base em estudos e transformações da sociedade, a definição de família também sofreu alterações.
Hoje a família nuclear, formada por pai, mãe e filho, divide lugar com outras formas de organização familiar, tais como : dois pais, duas mães, somente o pai, somente a mãe, avós, com tios, dentre outras formas.
Se durante muito tempo a família era conceituada por seus \textbf{membros} possuírem em comum \textbf{laços} sanguíneos, hoje o que lhe constitui é o respeito, o amor e principalmente, a \textbf{vontade} de estarem juntos, unidos e o importante que seja em prol do desenvolvimento integral de seus filhos.
Na direção de constituir uma relação famílias/escola pautada na parceria de um projeto educativo, FERREIRA (2017)⁷ sugere alguns verbos que concretizam objetivos a serem promovidos e geridos no projeto pedagógico : - \textbf{acolher} \textbf{bebês}, \textbf{crianças} e \textbf{adultos}, constituindo uma comunidade escolar ; - partilhar a construção permanente do projeto educativo de \textbf{bebês} e \textbf{crianças} \textbf{pequenas} ; e, - promover a participação, incluindo aqui os processos de tomada de decisões importantes para o projeto pedagógico da Educação \textbf{Infantil}.
A Diversidade na \textbf{Instituição} de Educação \textbf{Infantil} e a relação com a Família A \textbf{instituição} de Educação \textbf{Infantil} precisa construir caminhos para \textbf{acolher} bem todas as \textbf{crianças} deficientes, quilombolas, ribeirinhas, urbanas, rurais, extrativistas e florestais, conhecendo -as e respeitando -as como seres singulares que fazem parte de uma coletividade e que têm suas especificidades.
Para que isso seja possível, é imprescindível que todos os educadores percebam as famílias como \textbf{parceiras}, visto que, cada uma delas é o primeiro grupo social do qual a \textbf{criança} participa.
Assim, as famílias podem contribuir com informações e decisões importantes, colaborando para o planejamento da ação pedagógica.
As \textbf{instituições} de Educação \textbf{Infantil}, por sua vez, devem valorizar a diversidade cultural vivida pela \textbf{criança} no contexto familiar, e possibilitar a ampliação desse repertório. ⁷ Disponível em https://revistas.unilasalle.edu.br/index.php/Educacao/article/view/3646/pdf.
As famílias precisam ter garantidos o direito e o dever de compartilhar e colaborar com o que expressa o Art. 8º das DNCEI ( 2009 ), que dizem : > As propostas pedagógicas da Educação \textbf{Infantil} das \textbf{crianças} filhas de agricultores familiares, extrativistas, pescadores artesanais, ribeirinhos, assentados e acampados da reforma agrária, quilombolas, caiçaras, povos da floresta, devem : > > I - reconhecer os modos próprios de vida no campo como fundamentais para a constituição da identidade das \textbf{crianças} moradoras em territórios rurais ; > > II - ter vinculação inerente à realidade dessas populações, suas culturas, tradições e identidades, assim como as práticas ambientalmente sustentáveis ; > > III - flexibilizar, se necessário, \textbf{calendário}, \textbf{rotinas} e atividades respeitando as diferenças quanto à atividade econômica dessas populações ; > > IV - valorizar e evidenciar os saberes e o papel dessas populações na produção de conhecimentos sobre o mundo e sobre o ambiente natural ; > > V - prever a oferta de \textbf{brinquedos} e equipamentos que respeitem as características ambientais e socioculturais da comunidade ( DCNEI, 2009, parágrafo 3º ).
Diversas possibilidades devem ser disponibilizadas nas escolas de Educação \textbf{Infantil} para minimizar as ansiedades frente aos obstáculos que ainda enfrentam as famílias relacionadas às \textbf{crianças} com necessidades especiais.
A promoção de reuniões com associações de pais, professoras/professores, especialistas ou a estimulação para a criação de uma rede de apoio, de troca de experiências e de \textbf{ajuda} mútua são tarefas de todos, especialmente do grupo gestor de cada unidade escolar e do poder público, nas instâncias que o representam.
O trabalho desenvolvido com todas as \textbf{crianças} deve contribuir para que elas possam ser agentes transformadoras do seu contexto familiar, assim, é de suma importância que elas vivenciem, no cotidiano, experiências dotadas de significados, como por exemplo, a relação que as/os profissionais da \textbf{instituição} têm com as famílias fará parte da sua formação ética dos \textbf{meninos} e das \textbf{meninas}.
Nesse sentido, orientamos que as famílias sejam tratadas com respeito e igualdade, independentemente de sua raça, credo, poder aquisitivo, composição familiar, aparência física, local onde reside ou trabalha.
É necessário que as professoras, os professores tenham clareza de que nenhuma forma de preconceito é permitida dentro do espaço escolar.
Todos os momentos pedagógicos são antecedidos de planejamento, com a adaptação da \textbf{criança} não é diferente.
É comum \textbf{escutarmos} considerações de que a \textbf{criança} não está se adaptando.
Isso gera uma reflexão : é ela a responsável pelo sucesso ou insucesso desse período?
Assim, percebemos a responsabilidade dos \textbf{adultos}, a importância de repensar os momentos de acolhimento, possibilitando o envolvimento e a participação da família, de forma que sejam inseridas nessa rica experiência.
Ressaltamos ainda, que as \textbf{crianças} são diferentes.
Dessa forma, algumas adaptam -se ao novo ambiente, professora/professor e \textbf{colegas} com facilidade, enquanto outras demoram um tempo maior.
A \textbf{instituição} deve respeitar o ritmo e o tempo de cada \textbf{criança}, planejando estratégias acolhedoras e diferenciadas que atendam às suas individualidades.
Um diálogo permanente com a família é necessário, para que possam \textbf{conversar} sobre aspectos que influenciam no desenvolvimento e aprendizagem da \textbf{criança}, como sua história de vida, preferências e sentimentos ; a fim de traçar estratégias comuns para que ela se sinta apoiada, segura e motivada para realizar experiências individuais e coletivas, nos contextos, familiar e escolar, que contribuirão para o seu desenvolvimento integral.
Orienta -se que gestores, funcionários, professoras/professores e demais profissionais que lidam com a Educação \textbf{Infantil}, desenvolvam um trabalho com as famílias tendo como base, pontos relevantes discutidos nas Orientações Curriculares para a Educação \textbf{Infantil} ( 2011, p. 24 e 25 ), a seguir : - adotem uma atitude acolhedora em relação às famílias, de todas as \textbf{crianças}, sem discriminá -las por sua forma de organização : pais adolescentes, pais ou mães solteiras ( os ) ou divorciadas ( os ), casais homossexuais e outras ; - programem formas de \textbf{conversar} com as famílias, individualmente, ou em \textbf{pequenos} grupos, de modo a \textbf{acolher} suas expectativas, preocupações e trocar informações sobre as \textbf{crianças} ; - discutam com os pais o cotidiano e a proposta pedagógica da \textbf{Instituição} de Educação \textbf{Infantil}, por meio de \textbf{fotos}, de projeções de slides ou filmes, de exposições de produções \textbf{infantis}, de reuniões, ou da participação direta deles em algumas das atividades propostas às \textbf{crianças} ; - relatem aos pais algum episódio no qual a \textbf{criança} teve uma atuação espontânea, divertida, inteligente, habilidosa ; - \textbf{convidem} mães e pais ou outros familiares para vir tocar violão ou outros instrumentos musicais, \textbf{contar} histórias, ou produzir algo com as \textbf{crianças} ; - envolvam os pais em projetos tais como contação e/ou leitura de histórias para os filhos em casa, pesquisa sobre músicas regionais etc. ; - incentivem a participação dos pais nos conselhos escolares.
Nessa perspectiva, a família e a escola devem cultivar uma relação de confiança, na qual os saberes familiares sejam considerados, tais como : - incentivar a participação das famílias no processo de (re)elaboração da proposta pedagógica da \textbf{Instituição} ; - \textbf{acolher} e analisar as \textbf{sugestões} de mudança trazidas pelas famílias nessa construção, como também sensibilizando -as sobre a importância de conhecerem a \textbf{rotina}, compartilhando documentos específicos acerca do trabalho realizado na \textbf{instituição} junto às \textbf{crianças} e sobre os processos de desenvolvimento e aprendizagem ; - participarem das festividades, dos momentos de adaptação da \textbf{criança} ; - conhecerem de projetos desenvolvidos na \textbf{Instituição} ; - participarem de \textbf{passeios}, exposições, etc.
O Papel da Família no processo de transição da \textbf{criança} na/da Educação \textbf{Infantil} Quando se fala da transição aliada à vida das \textbf{crianças} com a família e a entrada na escola de Educação \textbf{Infantil} e no Ensino Fundamental, afirma -se que há um respaldo legal na própria Lei de Diretrizes e Bases nº 9394/96, art. 2º, no qual afirma que “ a educação é dever da família e do Estado e é importante que os pais garantam a educação escolar de seus filhos ”.
Brandão ( 2010 ) explica que esta relação deva ser inspirada nos princípios de liberdade proporcionando as condições necessárias para que a \textbf{criança} usufrua de seus direitos e dos seus ideais de solidariedade humana, tendo por finalidade seu desenvolvimento e qualificação.
E que na educação escolar, se for de qualidade, as \textbf{crianças} poderão obter o pleno desenvolvimento como educandos, devendo -se considerar que o papel dos pais e das \textbf{instituições} é formar cidadãos capazes de trabalhar e conviver harmonicamente em sociedade.
De acordo com as Orientações Curriculares para a Educação \textbf{Infantil} ( SEDUC, 2011 ), a perspectiva do atendimento aos direitos da \textbf{criança} na sua integralidade requer que as \textbf{instituições} de Educação \textbf{Infantil}, na organização de sua proposta pedagógica e curricular, assegurem espaços e tempos para participação, o diálogo e a \textbf{escuta} cotidiana das famílias, o respeito e a valorização das diferentes formas em que elas se organizam.
Isso deverá também e prioritariamente acontecer quando a \textbf{criança} sai da Educação \textbf{Infantil} para o Ensino Fundamental.
O Parecer nº 20/2009 do CNE \textbf{pontua} que : “ Quando a \textbf{criança} passa a frequentar a Educação \textbf{Infantil}, é preciso refletir sobre a especificidade de cada contexto no desenvolvimento da \textbf{criança} e a forma de integrar as ações e projetos educacionais das famílias e das \textbf{instituições} ”.
Nessa perspectiva, o trabalho pedagógico desenvolvido na Educação \textbf{Infantil} precisa considerar as diferentes formas de organização familiar, bem como seus saberes, fazeres e valores.
A atenção deverá ser redobrada quando acontecer a transição desta \textbf{criança} para o Ensino Fundamental.
Entende -se como necessário que a família conheça os objetivos da proposta escolar para acompanhar o desenvolvimento das práticas educativas das \textbf{crianças}, e se comprometa em acompanhar o sucesso na aprendizagem e na formação da \textbf{criança}.
A participação dos pais junto com os professores e demais educadores nos conselhos escolares e na organização do Projeto Político Pedagógico, nas festas e em outras atividades da escola possibilita agregar experiências e saberes e articular os dois contextos de desenvolvimento da \textbf{criança}.
Eventuais preocupações dos professores sobre a forma como algumas \textbf{crianças} parecem ser tratadas em casa — descuido, violência, discriminação, superproteção e outras — devem ser discutidas com o grupo gestor ou a coordenação pedagógica de cada unidade para que formas produtivas de esclarecimento e encaminhamento possam ser pensadas para beneficiar a família e a \textbf{criança}.
Por outro lado, compreende -se que a família também precisa ser conhecida e valorizada no contexto escolar, buscando -se sua integração e envolvimento.

% matched lemmas: acolher, adulto, ajuda, bebê, brinquedo, calendário, colega, contar, conversar, convidar, creche, criança, cuidado, cuidar, escuta, escutar, foto, infantil, infância, instituição, laço, membro, menino, parceiro, passeio, pequeno, pontuar, pré-escola, rotina, sugestão, vontade
\end{textsample}
